\subsectionaddtolist{الف}

سیستم فرکانس $\omega = \frac{\pi}{2}$ را به طور کامل حذف می‌کند. از آنجا که پاسخ ضربه حقیقی است، صفرهای سیستم باید به صورت جفت مزدوج مختلط روی دایره واحد در فرکانس‌های $\pm \frac{\pi}{2}$ قرار گیرند:
    \[
    z_{1,2} = e^{\pm j \frac{\pi}{2}} = \pm j
    \]
    بنابراین عامل صفرها در تابع تبدیل به صورت زیر است:
    \[
    (z - j)(z + j) = z^2 + 1
    \]

سیستم دارای دو قطب در شعاع $r = 0.5$ و زوایای $\theta = \pm \frac{\pi}{3}$ است:
    \[
    p_{1,2} = 0.5 e^{\pm j \frac{\pi}{3}} = 0.5 \left( \cos\frac{\pi}{3} \pm j \sin\frac{\pi}{3} \right) = 0.5 \left( \frac{1}{2} \pm j \frac{\sqrt{3}}{2} \right) = \frac{1}{4} \pm j \frac{\sqrt{3}}{4}
    \]
    بنابراین عامل قطب‌ها در مخرج تابع تبدیل برابر است با:
    \[
    (z - p_1)(z - p_2) = z^2 - (p_1 + p_2)z + p_1 p_2 = z^2 - 2(0.5)\cos\left(\frac{\pi}{3}\right)z + (0.5)^2
    \]
    \[
    = z^2 - 0.5z + 0.25
    \]

بنابراین شکل کلی تابع تبدیل با در نظر گرفتن ضریب بهره $K$ به صورت زیر خواهد بود:
\[
H(z) = K \frac{z^2 + 1}{z^2 - 0.5z + 0.25} = K \frac{1 + z^{-2}}{1 - 0.5z^{-1} + 0.25z^{-2}}
\]

برای یافتن $K$، از اطلاعات سوم یعنی بهره DC (در فرکانس $\omega = 0$ یا همان $z = e^{j0} = 1$) استفاده می‌کنیم که برابر با ۱ است:
\[
|H(1)| = 1 \implies K \frac{1^2 + 1}{1^2 - 0.5(1) + 0.25} = 1 \implies K \frac{2}{0.75} = 1 \implies K = \frac{0.75}{2} = \frac{3}{8} = 0.375
\]

بنابراین تابع تبدیل نهایی سیستم عبارت است از:
\[
H(z) = \frac{3}{8} \left( \frac{z^2 + 1}{z^2 - 0.5z + 0.25} \right) = 0.375 \left( \frac{1 + z^{-2}}{1 - 0.5z^{-1} + 0.25z^{-2}} \right)
\]

چون در صورت سوال ذکر شده که سیستم {علی} است، ناحیه همگرایی باید به صورت خارج بزرگترین قطب سیستم در صفحه Z باشد. بزرگترین شعاع قطب‌ها برابر با $r = 0.5$ است، پس:
\[
\text{ROC: } |z| > 0.5
\]

\subsectionaddtolist{ب}
بالاترین فرکانس نوسانی در سیستم‌های زمان-گسسته فرکانس نایکوئیست یعنی $\omega = \pi$ است که معادل جایگذاری $z = e^{j\pi} = -1$ در تابع تبدیل است:
\[
H(e^{j\pi}) = H(-1) = \frac{3}{8} \left( \frac{(-1)^2 + 1}{(-1)^2 - 0.5(-1) + 0.25} \right)
\]
\[
H(-1) = \frac{3}{8} \left( \frac{1 + 1}{1 + 0.5 + 0.25} \right) = \frac{3}{8} \left( \frac{2}{1.75} \right) = \frac{3}{8} \left( \frac{2}{\frac{7}{4}} \right) = \frac{3}{8} \times \frac{8}{7} = \frac{3}{7} \approx 0.428
\]

\subsectionaddtolist{ج}

یک سیستم IIR (که دارای قطب‌های غیر از مبدأ است) تنها زمانی می‌تواند دارای فاز خطی دقیق باشد که به ازای هر قطب در داخل دایره واحد، یک قطب متناظر در خارج دایره واحد به صورت معکوس یعنی $\frac{1}{p^*}$ داشته باشد تا تقارن لازم برقرار شود. 

اما در این سیستم:
\begin{enumerate}
    \item سیستم طبق فرض صورت سوال {علی} است.
    \item اگر بخواهیم قطب‌های معکوس در خارج دایره واحد (در شعاع $r=2$) اضافه کنیم تا فاز خطی شود، به دلیل علی بودن سیستم، ناحیه همگرایی مجبور است به سمتی فراتر از بزرگترین قطب (یعنی $|z| > 2$) برود که در این صورت دایره واحد را شامل نخواهد شد و سیستم ناپایدار می‌شود.
    \item از طرفی تابع تبدیل فعلی سیستم یک سیستم IIR مرتبه ۲ با دو قطب درون دایره واحد است و هیچ قطب متناظری در خارج دایره واحد برای خنثی‌سازی و ایجاد تقارن در پاسخ ضربه ندارد.
\end{enumerate}

بنابر دلایل فوق، این سیستم نمی‌تواند دارای فاز خطی باشد. سیستم‌های منتهی به فاز خطی دقیق با ساختار علی و پایدار صرفا در سیستم‌های FIR (که تمام قطب‌های آن‌ها در مبدأ قرار دارد) امکان‌پذیر است.
